% Standalone problem document, generated from the adjacent README.md.
% Compile with XeLaTeX. Mathematical content is maintained in Markdown.
\documentclass[11pt,a4paper]{article}
\usepackage[fontsize=10.5pt]{scrextend}
\usepackage[margin=22mm,headheight=15pt,headsep=9mm,footskip=11mm]{geometry}
\usepackage{fontspec}
\setmainfont{texgyrepagella}[Extension=.otf,UprightFont=*-regular,BoldFont=*-bold,ItalicFont=*-italic,BoldItalicFont=*-bolditalic]
\setsansfont{texgyreheros}[Extension=.otf,UprightFont=*-regular,BoldFont=*-bold,ItalicFont=*-italic,BoldItalicFont=*-bolditalic]
\setmonofont{texgyrecursor}[Extension=.otf,UprightFont=*-regular,BoldFont=*-bold,ItalicFont=*-italic,BoldItalicFont=*-bolditalic,Scale=MatchLowercase]
\usepackage{amsmath,amssymb}
\usepackage{unicode-math}
\setmathfont{texgyrepagella-math.otf}
\usepackage{xcolor,microtype,enumitem,needspace,fancyhdr}
\usepackage{bookmark}
\definecolor{ink}{HTML}{17344B}
\definecolor{accent}{HTML}{197D80}
\definecolor{muted}{HTML}{596773}
\hypersetup{colorlinks=true,linkcolor=accent,urlcolor=accent,pdftitle={TR-24 - Degree-five,
six, and nine generation of the Salmon tensor
ideal},pdfauthor={Open Problems in Numerical Linear Algebra}}
\urlstyle{same}
\setlength{\parindent}{0pt}
\setlength{\parskip}{5pt plus 1pt minus 1pt}
\linespread{1.04}
\setlength{\emergencystretch}{3em}
\widowpenalty=10000
\clubpenalty=10000
\displaywidowpenalty=10000
\allowdisplaybreaks[1]
\setlist{nosep,leftmargin=1.5em}
\providecommand{\tightlist}{\setlength{\itemsep}{0pt}\setlength{\parskip}{0pt}}
\providecommand{\pandocbounded}[1]{#1}
\pagestyle{fancy}
\fancyhf{}
\fancyhead[L]{\sffamily\footnotesize\color{muted}OPEN PROBLEMS IN NUMERICAL LINEAR ALGEBRA}
\fancyhead[R]{\sffamily\bfseries\footnotesize\color{accent}TR-24}
\fancyfoot[L]{\parbox[b]{0.9\textwidth}{\sffamily\fontsize{8}{10}\selectfont\color{muted}Editorial ratings; a bounded search cannot certify that a problem remains unsolved.\\Literature check: 2026-09-10}}
\fancyfoot[R]{\sffamily\footnotesize\color{muted}\thepage}
\renewcommand{\headrulewidth}{0.3pt}
\renewcommand{\headrule}{\hbox to\headwidth{\color{accent}\leaders\hrule height \headrulewidth\hfill}}
\makeatletter
\renewcommand{\section}{\Needspace{5\baselineskip}\@startsection{section}{1}{0pt}{12pt plus 2pt minus 2pt}{4pt}{\sffamily\bfseries\large\color{ink}}}
\renewcommand{\subsection}{\Needspace{4\baselineskip}\@startsection{subsection}{2}{0pt}{9pt plus 1pt minus 1pt}{3pt}{\sffamily\bfseries\normalsize\color{ink}}}
\makeatother
\setcounter{secnumdepth}{0}
\begin{document}
{\sffamily\small\color{muted}Tensor computations\par}
\vspace{3pt}
{\raggedright\hyphenpenalty=10000\exhyphenpenalty=10000\sffamily\bfseries\fontsize{21}{24}\selectfont\color{ink}Degree-five,
six, and nine generation of the Salmon tensor ideal\par}
\vspace{7pt}
{\sffamily\small\textbf{Difficulty:} extreme\quad\textbf{Importance:} interesting
to specialist\par}
{\sffamily\small\bfseries\color{accent}Status: Open\par}
\vspace{4pt}
{\color{accent}\hrule height 0.8pt}
\vspace{6pt}
\textbf{Rating rationale:} The remaining ideal-generation assertion is a
longstanding barrier in a prominent fixed-format tensor problem. Its
precise scheme-theoretic content principally matters to specialists in
tensor equations and algebraic statistics.

\subsection{Problem statement}\label{problem-statement}

Let \(V=\mathbb C^4\otimes\mathbb C^4\otimes\mathbb C^4\), with
coordinates \(t_{ijk}\), and let \[
X=\overline{\left\{
\sum_{\ell=1}^{4}a_\ell\otimes b_\ell\otimes c_\ell:
a_\ell,b_\ell,c_\ell\in\mathbb C^4
\right\}}^{\,\mathrm{Zar}}\subset V.
\] Thus \(X\) is the affine cone over the fourth secant variety of the
Segre embedding of \(\mathbb P^3\times\mathbb P^3\times\mathbb P^3\).
Its points are exactly the tensors of complex border rank at most four.

Let \[
I=I(X)=\{f\in\mathbb C[t_{ijk}]: f(T)=0\text{ for every }T\in X\}
\] be its homogeneous prime ideal. For each integer \(a\geq0\), let
\(I_a\) be the vector space of homogeneous degree-\(a\) polynomials in
\(I\). Is \[
I=\langle I_5\cup I_6\cup I_9\rangle
\] as an ideal of \(\mathbb C[t_{ijk}]\)?

The angle brackets mean all finite polynomial combinations of elements
of these three homogeneous spaces. This is the revised Salmon conjecture
in the degree-generation form explicitly stated by Friedland--Gross, §1.
It asks for equality of ideals, not merely equality of zero sets or
equality after taking radicals. It does not prescribe an unverified
number of generators or identify the full homogeneous pieces with
particular proposed representation modules.

\subsection{Why it matters}\label{why-it-matters}

The question asks for the tensor counterpart of determinantal equations
for low-rank matrices. Exact generators describe algebraic constraints
in small tensor decompositions and statistical mixture models, and
support symbolic and numerical methods for testing and studying these
constraints.

\newpage
\subsection{References and status
check}\label{references-and-status-check}

\begin{enumerate}
\def\labelenumi{\arabic{enumi}.}
\tightlist
\item
  S. Friedland and E. Gross, \emph{A proof of the set-theoretic version
  of the Salmon conjecture},
  \href{https://arxiv.org/pdf/1104.1776v2}{arXiv:1104.1776v2}, §1, p.1,
  states the revised degree-generation conjecture; Journal of Algebra
  \textbf{356} (2012), 374--379.
\item
  D. J. Bates and L. Oeding, \emph{Toward a salmon conjecture},
  \href{https://arxiv.org/pdf/1009.6181v2}{arXiv:1009.6181v2},
  introduction pp.1--2 and Theorem 3.10; Experimental Mathematics
  \textbf{20} (2011), 358--370.
\item
  J. Jagiełła and J. Jelisiejew, \emph{Unrestrictions and concise secant
  varieties},
  \href{https://arxiv.org/html/2604.24879v2}{arXiv:2604.24879v2},
  §1.5.3, explicitly distinguishes the solved set-theoretic problem from
  the unknown ideal.
\end{enumerate}

\subsubsection{Status check ---
2026-09-10}\label{status-check-2026-09-10}

Checked Friedland--Gross §1 against the displayed ideal equality and
Jagiełła--Jelisiejew v2, §1.5.3, together with targeted 2025--2026
Salmon-ideal searches. The 2026 source explicitly retains the ideal
question. The proved equality of zero sets only identifies the radical
of the proposed ideal; it does not settle a portion of this fixed
ideal-equality assertion. Numerical or local vanishing and dimension
certificates likewise do not prove ideal equality. No complete proof was
located, so the status remains Open.
\end{document}
